% GPT-6 Astra Pro; effort not exposed; /root/replacement; 22 November 2026.
\section[HD-F1-R31: the pure surface cancellation]{HD-F1-R31: the common quotient forces positive pure surface components and an actual defect extension}
\label{sec:followup-rm-pure-surface-cancellation}

Retain an actual object in the common-quotient branch
of R21--R30, with all its original hypotheses:
the genuine real-quadratic abelian fourfold, the
specified coherent faithful finite point action
and actual lift, the weakly semiregular gluable
adjacent triangle, and its equal-slope semistable
TF endpoints of ranks $(3,2)$. The actual
determinant twist has already made
$\operatorname{ch}(E)=v_q=(1-qA_1)(1-qA_2)$,
where $q>0$ is rational. The semistability
polarization is the original arbitrary ample $H$;
the forms $A_i=\theta_i^2/2$ use the retained
compatible RM polarization.

We prove three further restrictions on every
such actual object. Both pure RM coefficients
of the target's effective surface defect are
strictly positive. The connecting class becomes
zero after projection to the actual line hull.
Before that projection, its nonzero image in
the rank-one quotient is a boundary from the
actual hull defect of that quotient. The proof
also strengthens R22's required image from
at least three dimensions to a specified
four-dimensional subspace. These statements
retain, rather than remove, the remaining
surface-supported extension problem.

\paragraph{1. Fixing all actual sheaves and maps.}
Write the actual triangle and common quotient as
\[
 A\longrightarrow E\longrightarrow T[-1]
       \xrightarrow{\delta[-1]}A[1],\qquad
 \delta:T\longrightarrow A[2],\qquad
 p:A\twoheadrightarrow Q,\quad l:Q\hookrightarrow L=Q^{**}.
 \tag{F1R31.1}
\]
R22 gives invariant maps $u_1,u_2:Q\to T$
whose composites $u_bp$ form a basis of
$W=\operatorname{Hom}_G(A,T)$, together with
the actual equivariant hull isomorphism
$T^{**}=L^{\oplus2}$. For its inclusion and
quotient write
\[
 0\longrightarrow T\xrightarrow{j}L^{\oplus2}
       \xrightarrow{q_Z}Z\longrightarrow0,
 \qquad j=(j_1,j_2),\qquad j_a u_b=\delta_{ab}l.
 \tag{F1R31.2}
\]
The last equality holds first on the common
locally free open set by the defining hull
isomorphism, and then globally because $L$
is torsion free. Put
\[
 J=L/Q,\qquad
 \beta_Q=p[2]\delta\in\operatorname{Ext}^2_G(T,Q),
 \qquad \beta_L=l[2]\beta_Q\in\operatorname{Ext}^2_G(T,L).
 \tag{F1R31.3}
\]
All quotients here carry the induced actual
coherent lifts. No different extension class
or numerical replacement for a sheaf is used.

Use R26--R30's surviving whole characters
\[
 \operatorname{ch}(T)=a_0v_q+c\,p_1p_2,\qquad
 \operatorname{ch}(A)=(a_0+1)v_q+c\,p_1p_2,
 \quad
 p_1=r+\beta+sA_1,\quad p_2=\overline{p_1}.
 \tag{F1R31.4}
\]
Here $r,s\in F=\mathbf Q(\sqrt{\Delta_0})$,
$c\in\mathbf Q^*$, and the bar is the exact
factor-exchanging scalar-extension involution
of R26. The two nonzero divisor factors
$\beta,\overline\beta$ are definite by R30.
In particular $p_i$ is not a multiple of
$v_i=1-qA_i$. Let $\ell=c_1(L)=\ell_1+\ell_2$.
The rank and first-Chern identities are
\[
       a_0+c r\overline r=2,\qquad
       \ell_1=\frac{c\overline r}{2}\beta,
       \qquad \ell_2=\frac{cr}{2}\overline\beta.
 \tag{F1R31.5}
\]
For the actual nonzero effective surface cycle
$D=[Z]_2=\ell^2-\operatorname{ch}_2(T)$ write
\[
 \begin{split}
 D&=d_1A_1+d_2A_2+2\ell_1\ell_2-c\beta\overline\beta,\\
 d_1&=\frac{\ell_1^2}{A_1}+a_0q-c\overline r s,
       \qquad d_2=\overline{d_1}.
 \end{split}
 \tag{F1R31.6}
\]
The quotient in this display means its
coefficient in the one-dimensional space
$F A_1$. It is defined over $F$ by R26's
explicit projectors. Since the two pure
forms $A_i$ are semipositive and
$\int_X A_1A_2>0$,
\[
 d_1=\frac{\int_X D A_2}{\int_X A_1A_2}\ge0,
 \qquad
 d_2=\frac{\int_X D A_1}{\int_X A_1A_2}\ge0.
 \tag{F1R31.7}
\]
The mixed terms in (F1R31.6) wedge to zero
with either $A_i$ by their block degrees.
Field conjugacy gives $d_1=0$ exactly when
$d_2=0$. We shall exclude their simultaneous
vanishing using the actual support of $Z$.

\paragraph{2. The same coefficient occurs in both endpoint Yoneda maps.}
Let $K=K_1\oplus K_2$ be the balanced kernel
of R11 and R26. The two endpoint Hom
sequences give
\[
 \operatorname{ev}_A^2(K)\subset
      \{\delta f:f\in W\},\qquad
 \operatorname{ev}_T^2(K)\subset
      \{f[2]\delta:f\in W\}.
 \tag{F1R31.8}
\]
These are precisely R19's actual kernel
memberships. Define $\nu(f)=\delta f$
and $\mu(f)=f[2]\delta$. The traced
characteristic image of each endpoint
on $K$ has rank two by R26. As
$\dim W=2$, both maps
$\sigma_A^2\nu$ and $\sigma_T^2\mu$
are injective. Full Atiyah-trace cyclicity
gives the equality, for every actual $f$,
\[
                  \sigma_A^2(\delta f)
                         =\sigma_T^2(f[2]\delta).
 \tag{F1R31.9}
\]
Indeed naturality moves each Atiyah power
across $f$ and $\delta$, and graded trace
moves their composition cyclically. Their
degrees are zero and two, so the resulting
interchange sign is positive. This checks
every component of the full trace.
Since $\operatorname{ch}(A)-\operatorname{ch}(T)=v_q$
and $D_kv_q=0$ for $k\in K$, the two
traced characteristic images are equal.
The injectivity just proved therefore makes
the coefficients in (F1R31.8) identical,
unique and linear:
\[
 \begin{split}
 &f_K:K\twoheadrightarrow W,\qquad
       f_K(k)=u_{v(k)}p,\qquad
       u_v=v_1u_1+v_2u_2,\\
 &\operatorname{ev}_A^2(k)=\delta f_K(k),\qquad
       \operatorname{ev}_T^2(k)=u_{v(k)}[2]\beta_Q,
       \qquad v:K\twoheadrightarrow\mathbf C^2.
 \end{split}
 \tag{F1R31.10}
\]
Surjectivity follows from the same rank two.

Put $\mathcal H=H^1(X,\mathcal O_X)
=\mathcal H_1\oplus\mathcal H_2$, where each
RM block has dimension two. Choose nonzero
$\omega_i\in\bigwedge^2\mathcal H_i$ and write
\[
                         U=\mathbf C\omega_1\oplus\mathbf C\omega_2
                               \subset H^2(X,\mathcal O_X).
 \tag{F1R31.11}
\]
The block calculation (F1R26.11)--(F1R26.13)
defines nonzero linear functionals $f_i$ on
$K_i$ by $D_kp_i=f_i(k)\omega_i$, after
fixing the same nonzero scalar in $\omega_i$.
The action on (F1R31.4) has rank one on
each block, in two distinct block degrees.
The injectivity of $\sigma_T^2\mu$ then gives
linearly independent columns $v_1',v_2'$
such that
\[
                         v(k)=f_i(k)v_i'\quad(k\in K_i).
 \tag{F1R31.12}
\]
All fixed nonzero constants, including $c$,
are absorbed in these columns.

Likewise the actual line characteristic map
has
\[
 \operatorname{ev}_L^2(k)=g_i(k)\omega_i\quad(k\in K_i),
          \qquad D_k e^{\ell_i}=g_i(k)\omega_i.
 \tag{F1R31.13}
\]
Multiplication by $e^{\ell_i}$ fixes $\omega_i$
because any positive divisor degree would
exceed antiholomorphic degree two in that
block; the other block factor cancels in
the line's injective full trace. The kernel
of the map from the six-dimensional block
Hodge space to $K_i^*$ is exactly
$\mathbf C v_i$, by the nondegenerate
pairing (F1R26.11). Thus $g_i\ne0$:
$g_i=0$ would imply $e^{\ell_i}=a v_i$;
degree zero gives $a=1$, degree two
gives $\ell_i=0$, and degree four then
gives $0=-qA_i$. Consequently
\[
                         \operatorname{ev}_L^2(K)=U.
 \tag{F1R31.14}
\]
Naturality on the actual rows of $j$ is
the exact identity
\[
     g_i(k)\omega_i j_b
           =f_i(k)(v_i')_b\beta_L
                   \quad(k\in K_i,\ b=1,2).
 \tag{F1R31.15}
\]

\paragraph{3. Surfaces with both pure restrictions zero.}
We prove the geometric fact needed for
the actual cycle in (F1R31.6). Let $S$
be an irreducible surface in $X$ whose
two pure holomorphic two-forms restrict
to zero on a resolution $g:Y\to S\subset X$.
At a general smooth point its tangent
plane has zero pure Pl\"ucker coefficients
in both two-dimensional eigenspaces
$T_0X=V_1\oplus V_2$. Each projection
therefore has rank at most one. Their
combined rank is two, so both have rank
one and the tangent plane is exactly
\[
                        T_yS=\ell_1(y)\oplus\ell_2(y),
                                   \quad\ell_i(y)\subset V_i.
 \tag{F1R31.16}
\]
In this paragraph the symbols $\ell_i(y)$
denote tangent lines, not divisor classes.

After a translation the actual period
construction of the Albanese gives
$g=\psi a_Y$, with
$\psi:\operatorname{Alb}(Y)\to X$ a
homomorphism. Let $B=\operatorname{im}\psi$.
Its cotangent kernel is exactly
$\ker(g^*:H^0(\Omega_X^1)\to H^0(\Omega_Y^1))$.
Put $E_i=g^*(V_i^*)$.
At the general points in (F1R31.16),
the two spaces evaluate in distinct
lines of $T_y^*Y$. A global form in
$E_1\cap E_2$ therefore vanishes on
this dense open set and is zero.
It follows that the cotangent kernel
splits blockwise. Its annihilator
$T_0B$ is preserved by the actual
endomorphism
\[
                  j_F=m(1+\sqrt{\Delta_0})\in\operatorname{End}(X),
 \tag{F1R31.17}
\]
whose distinct real eigenvalues are
$\alpha_1=m(1+\sqrt{\Delta_0})$ and
$\alpha_2=m(1-\sqrt{\Delta_0})$.
Thus $j_F(B)\subset B$: the image of a
connected subtorus is determined by
its tangent space. Each $E_i$ has
dimension one or two, and
$\dim B=\dim E_1+\dim E_2$.

There is no nonzero abelian subvariety
with tangent contained in one $V_i$.
Such a subvariety would be preserved by
$j_F$, whose restriction on its real
universal cover is scalar multiplication
by the real irrational number $\alpha_i$.
It would induce that same scalar map on
its nonzero rational first homology,
which is impossible for a rational
linear endomorphism: its matrix in a
rational basis would be $\alpha_i I$.
The identical argument excludes
$\dim B=3$, because the actual quotient
$X\to X/B$ would be an elliptic quotient
on which $j_F$ induces one of these
irrational real scalars.

It remains to exclude $\dim B=4$.
Then each $E_i$ is two dimensional and
its two independent holomorphic forms
have wedge zero. The published
Castelnuovo--de Franchis theorem gives
actual morphisms with connected fibers
\[
                 f_i:Y\longrightarrow C_i,
                 \qquad g(C_i)\ge2,\qquad
                 E_i\subset f_i^*H^0(C_i,\Omega_{C_i}^1).
 \tag{F1R31.18}
\]
The two different cotangent directions
in (F1R31.16) make $(f_1,f_2)$ generically
of rank two. It is consequently a
surjective generically finite morphism
onto $C_1\times C_2$. Its actual
Albanese homomorphism
\[
          a:\operatorname{Alb}(Y)\twoheadrightarrow
                            JC_1\times JC_2
 \tag{F1R31.19}
\]
is surjective: pullback of holomorphic
one-forms along a dominant generically
finite map is injective, and its dual
is the differential of $a$.

Now $\psi^*H^0(\Omega_X^1)=E_1\oplus E_2$
is contained in the image of $a^*$.
The homomorphism $\psi$ kills
$(\ker a)^0$, because its differential
vanishes there. Factor it through
$C=\operatorname{Alb}(Y)/(\ker a)^0$.
The induced map $\bar a:C\to JC_1\times JC_2$
is an isogeny. Choose an integer $M$
annihilating its finite kernel. Multiplication
by $M$ on $C$ descends to an actual
homomorphism $t:JC_1\times JC_2\to C$
with $\bar a t=[M]$ and $t\bar a=[M]$.
Writing $\bar\psi:C\to X$, the composite
$\gamma=\bar\psi t$ satisfies
\[
                              \gamma a=[M]\psi.
 \tag{F1R31.20}
\]
Injectivity of $a^*$ and (F1R31.18)
show that $\gamma^*(V_i^*)$ lies entirely
in the $JC_i$ cotangent summand. Thus
$\gamma(JC_i)$ has tangent contained in
$V_i$. Equation (F1R31.20) and $B=X$
make $\gamma$ surjective, so these pure
abelian subvarieties are nonzero. This
contradicts the preceding rational-homology
argument. Therefore $\dim B=2$.
Since $S$ is a closed irreducible surface
contained in a translate of the abelian
surface $B$, it is that translate.
Its tangent has one line in each block
and $B$ is preserved by $j_F$.

Any two distinct such abelian surfaces
have finite intersection. Indeed the
identity component of their intersection
is preserved by $j_F$. A one-dimensional
intersection has a $j_F$-stable tangent
line; because its eigenvalues are distinct,
that line lies in one $V_i$, already
excluded. A two-dimensional intersection
makes the two surfaces equal. The same
conclusion holds for translates, since a
nonempty intersection is a coset of the
intersection subgroup. Distinct translates
of one surface are disjoint.

\paragraph{4. The zero-pure-cycle case contradicts the actual cancellation image.}
Suppose $d_1=d_2=0$. By (F1R31.7) every
irreducible component of the effective
cycle $D$ has zero integral against
both $A_i$. On its resolution, $A_i$
is a positive constant times the square
norm of the restricted pure holomorphic
two-form. Thus both forms are zero,
and paragraph~3 applies to every
surface component.

Let $Z_0\subset Z$ be its maximal
subsheaf of dimension at most one and
put $P=Z/Z_0$. The precise vanishing
and actual comparison from R30 are
\[
 \operatorname{Ext}^1_X(Z_0,L)
  =\operatorname{Ext}^2_X(Z_0,L)=0,
 \qquad \operatorname{Ext}^2_X(P,L)
                   \xrightarrow{\sim}\operatorname{Ext}^2_X(Z,L).
 \tag{F1R31.21}
\]
The pure sheaf $P$ is supported on
finitely many distinct translates $S_a$
of the abelian surfaces just proved.
Their pairwise intersections are finite.
This allows an actual decomposition of
the relevant Ext group even when the
surfaces have different tangent spaces.

For each $a$ let
$V_a=X\setminus\bigcup_{b\ne a}S_b$ and
let $P_a$ be the image of
$P\to (\iota_a)_*(P|_{V_a})$.
It is a coherent quotient of $P$:
its kernel is a coherent subsheaf of
the Noetherian sheaf $P$. Its support
is $S_a$ and it is pure of dimension
two. To check purity, a lower-dimensional
subsheaf restricts to zero on $V_a$
by purity of $P|_{V_a}$, and is then
zero by the displayed injection into
$(\iota_a)_*(P|_{V_a})$.
The resulting actual map gives
\[
             0\longrightarrow P\longrightarrow
                    \bigoplus_a P_a\longrightarrow C_0
                           \longrightarrow0,
                    \qquad \dim\operatorname{Supp}(C_0)\le0.
 \tag{F1R31.22}
\]
Its kernel is generically zero on
every surface and hence zero by
purity. Away from the finite pairwise
intersections it is an isomorphism,
which proves the assertion about
$C_0$. Each $P_a$ keeps the full
generic infinitesimal thickening along
$S_a$; no restriction to reduced
support has been made.
Codimension four gives
$\operatorname{Ext}^2_X(C_0,L)=
\operatorname{Ext}^3_X(C_0,L)=0$.
The two adjacent terms of the Hom
sequence of (F1R31.22) therefore give
the actual isomorphism
\[
       \bigoplus_a\operatorname{Ext}^2_X(P_a,L)
                   \xrightarrow{\sim}\operatorname{Ext}^2_X(P,L).
 \tag{F1R31.23}
\]
This comparison is made for ordinary
Ext groups and needs no equivariant
choice of the individual components.

Write $S_a=x_a+B_a$, and let
$\pi_a:X\to X/B_a$ be the actual
quotient homomorphism. Its base has
dimension two. The entire sheaf
$\mathcal F_a=P_a\otimes L^{-1}$ is
annihilated by a power of the pullback
of the ideal of $\pi_a(x_a)$. Choose
an affine open in $X/B_a$ containing
that point and its full finite
infinitesimal neighborhood. As proved
by the restriction and local-to-global
Ext maps in (F1R30.8)--(F1R30.14),
every class in
\[
         U_a=\pi_a^*H^1(X/B_a,\mathcal O_{X/B_a})
                       \subset\mathcal H
 \tag{F1R31.24}
\]
acts by the zero morphism
$\mathcal F_a\to\mathcal F_a[1]$.
Indeed its restriction on the affine
base is zero, while restriction of
self-Ext to its inverse image is an
isomorphism because the full support
is closed inside that inverse image.
The same proof retains all nilpotents.

The tangent space of $B_a$ has one
line in each RM block. Thus
$U_a=\ker(\mathcal H\to H^1(B_a,\mathcal O_{B_a}))$
has one line in each $\mathcal H_i$,
and
\[
                            U\subset U_a\wedge\mathcal H.
 \tag{F1R31.25}
\]
Let $q_a:L^{\oplus2}\to P_a$ be the
actual composite of $q_Z$ with the
maps in (F1R31.21)--(F1R31.22).
Tensor with $L^{-1}$. Naturality of
the cohomology action on each of
its two columns and (F1R31.24) imply
\[
        (q_a)_*:
             H^2(X,\mathcal O_X)^{\oplus2}
                    \longrightarrow H^2(X,\mathcal F_a)
                    \quad\text{kills }U\otimes\mathbf C^2.
 \tag{F1R31.26}
\]
For example a product $\alpha[1]\xi$
with $\alpha\in U_a$ is sent by a
column $s$ to
$\alpha_{\mathcal F_a}[1]s[1]\xi=0$.

Fix an invariant volume form on $X$.
CY4 Serre duality identifies the
transpose of (F1R31.26) with the
actual precomposition map
\[
       q_a^*:\operatorname{Ext}^2_X(P_a,L)
                     \longrightarrow
                     H^2(X,\mathcal O_X)\otimes(\mathbf C^2)^*.
 \tag{F1R31.27}
\]
The equality is the literal equality
of traces of the compositions
$b[2]q_a[2]c$, as in R30, so it
uses the same quotient $q_a$.
On $H^2(\mathcal O_X)$ the pairing
is wedge product into $H^4(\mathcal O_X)$.
It is nondegenerate on $U$, since
$\omega_1\omega_2\ne0$. Its orthogonal
complement is precisely the mixed
space $\mathcal H_1\otimes\mathcal H_2$.
Hence (F1R31.26)--(F1R31.27) give
\[
        \operatorname{im}q_a^*
             \subset U^\perp\otimes(\mathbf C^2)^*.
 \tag{F1R31.28}
\]
The isomorphisms (F1R31.21) and
(F1R31.23) identify $q_Z^*$ with
the sum of these actual maps. Thus
its ordinary image has zero
intersection with
$U\otimes(\mathbf C^2)^*$.
But (F1R31.14) and the precise
subspace statement (F1R22.16) require
\[
 \dim\left(\operatorname{im}q_Z^*
       \cap\bigl(U\otimes(\mathbf C^2)^*\bigr)\right)\ge3
 \tag{F1R31.29}
\]
already on invariants. This is a
contradiction. Consequently
\[
                                 d_1>0,\qquad d_2>0.
 \tag{F1R31.30}
\]

\paragraph{5. Positive pure coefficients force the precise defect extension.}
Inside the actual six-dimensional block
Hodge space, (F1R31.5)--(F1R31.6) give
the identity
\[
 e^{\ell_1}-\frac{a_0}{2}v_1
                   -\frac{c\overline r}{2}p_1
                       =\frac{d_1}{2}A_1,
 \tag{F1R31.31}
\]
and its conjugate in block two.
For clarity, its rank component is
$1-a_0/2-cr\overline r/2=0$, its
divisor component is zero by
(F1R31.5), and its top component
is exactly $d_1A_1/2$.

The pairing calculation used in
(F1R31.13) proves that $g_i$ and
$f_i$ are proportional exactly when
$e^{\ell_i}-b p_i$ lies in
$\mathbf C v_i$ for some $b$.
Since the divisor component of
$p_i$ is nonzero, it forces
$b=c\overline r/2$ in block one
and $b=cr/2$ in block two.
The rank component then forces the
coefficient of $v_i$ to be $a_0/2$.
Thus proportionality is equivalent
to $d_i=0$. This also covers $r=0$:
the required coefficient $b$ is then
zero, and (F1R31.31) still applies.
By (F1R31.30), the two functionals
are independent in each block.
Choose $k\in K_i$ with
$f_i(k)=0$ and $g_i(k)\ne0$.
The actual row identity (F1R31.15)
therefore gives
\[
                            \omega_i j_b=0
                                 \quad(i=1,2;\ b=1,2).
 \tag{F1R31.32}
\]
Returning to (F1R31.15), or to its
sum over both blocks, and using
surjectivity of $v$ in (F1R31.10),
we obtain
\[
                              \beta_L=l[2]p[2]\delta=0.
 \tag{F1R31.33}
\]
The equivariant Hom sequence of
(F1R31.2) now strengthens (F1R31.29)
to the exact required inclusion
\[
 U\otimes(\mathbf C^2)^*
       \subset\operatorname{im}\left[
       \operatorname{Ext}^2_G(Z,L)
               \xrightarrow{q_Z^*}
               \operatorname{Ext}^2_G(L^{\oplus2},L)\right].
 \tag{F1R31.34}
\]
This is a four-dimensional space.
Translations act trivially on
$H^2(\mathcal O_X)$, line tensors
cancel in the endomorphism action,
and the two invariant columns used
in the actual hull isomorphism make
its multiplicity space trivial.
Thus the entire displayed space is
invariant for the given lift, exactly
as checked in R22.

Finally apply equivariant Hom from
$T$ to $0\to Q\xrightarrow{l}L\to J\to0$.
The actual segment is
\[
 \operatorname{Ext}^1_G(T,J)
        \xrightarrow{\partial}
 \operatorname{Ext}^2_G(T,Q)
        \xrightarrow{l_*}
 \operatorname{Ext}^2_G(T,L).
 \tag{F1R31.35}
\]
Equation (F1R31.33) supplies an actual
class $\xi\in\operatorname{Ext}^1_G(T,J)$
such that
\[
                      0\ne\beta_Q=p[2]\delta=\partial\xi.
 \tag{F1R31.36}
\]
Its nonvanishing follows from
(F1R31.10): otherwise the target's
whole characteristic action on $K$
would be zero, contrary to R26's
rank two. Exactness on invariants
uses averaging for the finite
actual coherent action in
characteristic zero.

This does not lift $\delta$ itself
through the defect $J$. The remaining
source correction is explicit. For
$K_A=\ker p$, naturality on $p$ gives
\[
       (\operatorname{ev}_Q^2(k)
                   -\beta_Q u_{v(k)})p=0.
 \tag{F1R31.37}
\]
The difference belongs to the image
of the actual preceding map
$\operatorname{Ext}^1_G(K_A,Q)
\to\operatorname{Ext}^2_G(Q,Q)$
in the Hom sequence of
$0\to K_A\to A\to Q\to0$.
No vanishing of that group or
of this correction has been asserted.

\paragraph{Published framework, exact checking scope, and remaining gate.}
The two rank-two trace maps and their
block pairings are the explicit R19,
R26 calculations in the published
Buchweitz--Flenner characteristic/trace
framework. Equation (F1R31.9) uses
its actual naturality and cyclic trace,
as verified component by component above.
The only new geometric fibration input
is the classical Castelnuovo--de Franchis
theorem, in the $k=1$ case of
M.~Brunella, \emph{On holomorphic forms
on compact complex threefolds},
Comment.~Math.~Helv.~74 (1999), 642--656,
Lemma~1, pp.~645--646. Its complete
published proof was read. The full
proof of the classical statement and
its meromorphic-ratio extension step
were also read in P.~Py,
\emph{Lectures on K\"ahler Groups},
Princeton University Press, 2025,
using the author's 6 November 2024
text, Theorem~4.1 and Proposition~4.4,
pp.~46--49. This source comparison
does not identify the page numbering
of the author text with that of the
published volume.
The actual quotient-torus maps,
support excision, full infinitesimal
thickenings, Serre transpose and both
codimension vanishings are the retained
R30 maps, using the complete published
Huybrechts--Lehn Proposition~1.1.6(i)
and the previously audited support
and affine-cohomology results.

The new conclusions are precisely
(F1R31.30), (F1R31.33), (F1R31.34)
and (F1R31.36). The remaining common
quotient must have the indicated
positive pure surface coefficients,
the four-dimensional invariant
cancellation image, and a nonzero
actual defect extension coupled to
the original $\delta$ through
(F1R31.36)--(F1R31.37). This proof
constructs none of those remaining
objects and asserts no further
vanishing of their surface Ext groups.
